\documentclass[a5paper,10pt,dvipdfmx]{jsarticle}
\usepackage{amsmath,amssymb,amsthm}
\usepackage[margin=17mm,headsep=5mm,includehead]{geometry}
\usepackage{fancyhdr}
\pagestyle{fancy}
\fancyhf{}
\fancyhead[L]{\small 数学原論\quad N-7}
\fancyhead[R]{\small 濃度と可算性\quad\thepage}
\renewcommand{\headrulewidth}{0pt}
\setlength{\headheight}{14pt}
\theoremstyle{definition}
\newtheorem{ruledef}{規約}[section]
\newtheorem{principle}[ruledef]{原理}
\newtheorem{defn}[ruledef]{定義}
\theoremstyle{plain}
\newtheorem{prop}[ruledef]{命題}
\newtheorem{thm}[ruledef]{定理}
\renewcommand{\proofname}{証明}
\newcommand{\iffc}{\Longleftrightarrow}
\newcommand{\impc}{\Longrightarrow}
\usepackage{needspace}
\newcommand{\Pow}{\mathcal{P}}
\newcommand{\id}{\mathrm{id}}
\newcommand{\emptysetc}{\varnothing}
\clubpenalty=10000\widowpenalty=10000\displaywidowpenalty=10000

\begin{document}
\begin{center}
{\small 数学原論}\par\vspace{4mm}
{\LARGE 濃度と可算性}\par\vspace{2mm}
{\small N-7\quad 第1版}
\end{center}
本篇は N-5「自然数と有限集合」及びその前提を用いる。集合の濃度を写像による比較として導入し，Cantor--Bernstein の定理，Cantor の定理，可算集合の理論を展開する。選択公理及び一般の整列定理を用いない。
\Needspace{8\baselineskip}
\section{集合の比較}
% genron:begin K-000133
% genron:use K-000001
% genron:use K-000002
% genron:use K-000003
% genron:use K-000004
% genron:use K-000005
% genron:use K-000007
% genron:use K-000017
% genron:use K-000039
% genron:use K-000040
% genron:use K-000041
% genron:use K-000052
% genron:use K-000110
% genron:use K-000111
\Needspace{4\baselineskip}
\begin{defn}[濃度の比較]\label{k133}
集合 $A,B$ の間に全単射が存在するとき，両者は同じ濃度を持つといい $|A|=|B|$ と記す。単射 $A\to B$ が存在するとき $|A|\leq |B|$ と記す。一般の集合について $|A|$ を新たな集合として構成しているのではなく，これらの式全体を写像の存在の略記として用いる。有限集合では N-5 の元の個数と整合する（K-000111）。
\end{defn}
% genron:end K-000133
% genron:begin K-000134
% genron:use K-000001
% genron:use K-000002
% genron:use K-000003
% genron:use K-000004
% genron:use K-000005
% genron:use K-000007
% genron:use K-000017
% genron:use K-000036
% genron:use K-000039
% genron:use K-000040
% genron:use K-000041
% genron:use K-000133
\Needspace{4\baselineskip}
\begin{prop}[濃度の比較の性質]\label{k134}
同濃度の関係は反射的，対称的及び推移的である。単射による比較は反射的及び推移的で，同濃度の集合への置換によって変わらない。
\end{prop}
\begin{proof}
恒等写像，逆写像及び合成を用いれば直ちに従う。
\end{proof}
% genron:end K-000134
% genron:begin K-000135
% genron:use K-000001
% genron:use K-000002
% genron:use K-000003
% genron:use K-000004
% genron:use K-000005
% genron:use K-000007
% genron:use K-000017
% genron:use K-000024
% genron:use K-000027
% genron:use K-000029
% genron:use K-000035
% genron:use K-000036
% genron:use K-000037
% genron:use K-000038
% genron:use K-000039
% genron:use K-000040
% genron:use K-000041
% genron:use K-000042
% genron:use K-000045
% genron:use K-000097
% genron:use K-000102
% genron:use K-000105
% genron:use K-000107
% genron:use K-000133
\Needspace{4\baselineskip}
\begin{thm}[Cantor--Bernstein の定理]\label{k135}
集合 $A,B$ について $|A|\leq |B|$ かつ $|B|\leq |A|$ ならば $|A|=|B|$ である。
\end{thm}
\begin{proof}
単射 $f\colon A\to B$ と $g\colon B\to A$ を取る。$T=g\circ f$ とし，再帰定理を値域 $\Pow(A)$ に適用して
\[ A_0=A\setminus g(B),\qquad A_{n+1}=T(A_n) \]
と定める。$D=\bigcup_{n\in\mathbb N}A_n$ とおく。このとき $D=A_0\cup T(D)$ であり，$A_0$ と $T(D)$ は交わらない。実際，後者は $g(B)$ に含まれ，$\mathbb N$ の各元は $0$ 又は後者であるから合併の等式を得る。

$A\setminus D\subseteq g(B)$ である。$g$ の終域を $g(B)$ に縮小すれば全単射であるから，次の写像が定まる。
\[ h(a)=\begin{cases}f(a)&(a\in D),\\g^{-1}(a)&(a\notin D).\end{cases} \]
各部分上で $h$ は単射である。もし $a\in D$，$a'\notin D$ について $f(a)=g^{-1}(a')$ なら $a'=T(a)\in D$ となり矛盾する。よって $h$ は単射である。任意の $b\in B$ に対し，$b\in f(D)$ なら $b$ は $h$ の像に属する。$b\notin f(D)$ なら $g(b)\notin T(D)$（$g$ の単射性による）であり $g(b)\notin A_0$ であるから $g(b)\notin D$ となり $h(g(b))=b$ である。従って $h$ は全射でもある。
\end{proof}
% genron:end K-000135
% genron:begin K-000136
% genron:use K-000001
% genron:use K-000002
% genron:use K-000003
% genron:use K-000004
% genron:use K-000005
% genron:use K-000007
% genron:use K-000017
% genron:use K-000024
% genron:use K-000035
% genron:use K-000037
% genron:use K-000039
% genron:use K-000041
% genron:use K-000042
% genron:use K-000133
\Needspace{4\baselineskip}
\begin{prop}[全射と濃度の比較]\label{k136}
$A\ne\varnothing$ とし，単射 $f\colon A\to B$ が存在するとする。このとき全射 $B\to A$ が存在する。
\end{prop}
\begin{proof}
$a_0\in A$ を一つ取る。$f(A)$ 上では $f$ の逆写像，その外では定値 $a_0$ を取る写像 $g\colon B\to A$ を作れば $g\circ f=\id_A$ となる。従って $g$ は全射である。
\end{proof}
% genron:end K-000136
% genron:begin K-000137
% genron:use K-000001
% genron:use K-000002
% genron:use K-000003
% genron:use K-000004
% genron:use K-000005
% genron:use K-000007
% genron:use K-000017
% genron:use K-000023
% genron:use K-000024
% genron:use K-000027
% genron:use K-000035
% genron:use K-000039
% genron:use K-000041
% genron:use K-000133
\Needspace{4\baselineskip}
\begin{thm}[Cantor の定理]\label{k137}
任意の集合 $A$ に対し $|A|\leq|\Pow(A)|$ が成り立つが，全射 $A\to\Pow(A)$ は存在しない。従って両者は同濃度でない。$|A|\leq|B|$ かつ $|A|\ne|B|$ のとき $|A|<|B|$ と記すことにすれば，$|A|<|\Pow(A)|$ が成り立つ。
\end{thm}
\begin{proof}
$a\mapsto\{a\}$ は単射である。全射 $f\colon A\to\Pow(A)$ が存在すると仮定し，$D=\{a\in A\mid a\notin f(a)\}$ とおく。$D\in\Pow(A)$ であるから $f(d)=D$ となる $d\in A$ が存在する。このとき $d\in D\iffc d\notin D$ となり矛盾する。
\end{proof}
% genron:end K-000137
\Needspace{8\baselineskip}
\section{可算集合}
% genron:begin K-000138
% genron:use K-000001
% genron:use K-000002
% genron:use K-000003
% genron:use K-000004
% genron:use K-000005
% genron:use K-000007
% genron:use K-000017
% genron:use K-000039
% genron:use K-000040
% genron:use K-000099
% genron:use K-000101
% genron:use K-000110
% genron:use K-000133
\Needspace{4\baselineskip}
\begin{defn}[高々可算・可算無限]\label{k138}
$\mathbb N$ への単射を持つ集合を高々可算という。$\mathbb N$ と全単射を持つ集合を可算無限という。有限集合は自然数の $\mathbb N$ への包含と全単射との合成により，高々可算である。
\end{defn}
% genron:end K-000138
% genron:begin K-000139
% genron:use K-000001
% genron:use K-000002
% genron:use K-000003
% genron:use K-000004
% genron:use K-000005
% genron:use K-000007
% genron:use K-000017
% genron:use K-000039
% genron:use K-000040
% genron:use K-000097
% genron:use K-000099
% genron:use K-000101
% genron:use K-000103
% genron:use K-000110
% genron:use K-000111
% genron:use K-000138
\Needspace{4\baselineskip}
\begin{prop}[自然数集合の無限性]\label{k139}
$\mathbb N$ は有限集合でない。
\end{prop}
\begin{proof}
もし $\mathbb N$ と自然数 $n$ の間に全単射があれば，$S(n)\subseteq\mathbb N$ の包含と合成することによって単射 $S(n)\to n$ を得る。N-5 の有限集合の単射比較（K-000111）に反する。
\end{proof}
% genron:end K-000139
% genron:begin K-000140
% genron:use K-000001
% genron:use K-000002
% genron:use K-000003
% genron:use K-000004
% genron:use K-000005
% genron:use K-000007
% genron:use K-000017
% genron:use K-000024
% genron:use K-000035
% genron:use K-000039
% genron:use K-000068
% genron:use K-000076
% genron:use K-000099
% genron:use K-000100
% genron:use K-000103
% genron:use K-000104
% genron:use K-000105
% genron:use K-000111
% genron:use K-000113
% genron:use K-000139
\Needspace{4\baselineskip}
\begin{thm}[自然数の無限部分集合]\label{k140}
$B\subseteq\mathbb N$ が有限集合でなければ，$B$ と $\mathbb N$ は順序同型である。
\end{thm}
\begin{proof}
$B$ は空でない。任意の $b\in B$ に対し，$b$ より大きい $B$ の元が存在する。存在しないならば $B\subseteq S(b)$ となり有限である。よって $b$ より大きい $B$ の元のうち最小のものを $t(b)$ とする。最小元の一意性から写像 $t\colon B\to B$ が得られる。再帰定理により $u\colon\mathbb N\to B$ を，$u(0)=\min B$，$u(S(n))=t(u(n))$ によって定める。$u(n)<u(S(n))$ であるから，$m<n$ なら $u(m)<u(n)$ となることが $n$ に関する帰納法で従う。従って $u$ は順序埋込みである。

$b\in B\setminus u(\mathbb N)$ が存在すると仮定する。もし $u(n)>b$ となる $n$ があれば，その最小元を取る。$n=0$ は $u(0)$ の定義に反する。$n=S(k)$ なら $u(k)<b<u(S(k))$ であり，$t$ の定義に反する。従ってすべての $n$ で $u(n)<b$ となる。しかしこの場合，$u$ の $S(b)$ への制限は $b$ への単射を与え，有限集合の単射比較に反する。従って $u$ は全射であり，順序同型である。
\end{proof}
% genron:end K-000140
% genron:begin K-000141
% genron:use K-000001
% genron:use K-000002
% genron:use K-000003
% genron:use K-000004
% genron:use K-000005
% genron:use K-000007
% genron:use K-000017
% genron:use K-000035
% genron:use K-000037
% genron:use K-000039
% genron:use K-000041
% genron:use K-000042
% genron:use K-000104
% genron:use K-000110
% genron:use K-000113
% genron:use K-000138
% genron:use K-000139
% genron:use K-000140
\Needspace{4\baselineskip}
\begin{prop}[高々可算集合の分類]\label{k141}
高々可算集合は有限又は可算無限のいずれか一方である。高々可算集合の任意の部分集合及び写像による像も高々可算である。
\end{prop}
\begin{proof}
高々可算集合 $A$ は $\mathbb N$ の部分集合と全単射を持つので，定理\ref{k140}から前半が従う。部分集合は単射の制限によって高々可算である。像については $A$ を $B\subseteq\mathbb N$ に移し，$f\colon B\to X$ を考える。各 $x\in f(B)$ に対し $f^{-1}(\{x\})$ の最小元を対応させる写像は $f(B)\to\mathbb N$ の単射である。従って選択公理は不要である。
\end{proof}
% genron:end K-000141
% genron:begin K-000142
% genron:use K-000001
% genron:use K-000002
% genron:use K-000003
% genron:use K-000004
% genron:use K-000005
% genron:use K-000007
% genron:use K-000017
% genron:use K-000024
% genron:use K-000033
% genron:use K-000035
% genron:use K-000039
% genron:use K-000068
% genron:use K-000069
% genron:use K-000072
% genron:use K-000097
% genron:use K-000099
% genron:use K-000103
% genron:use K-000110
% genron:use K-000111
% genron:use K-000113
% genron:use K-000116
% genron:use K-000135
% genron:use K-000138
\Needspace{4\baselineskip}
\begin{thm}[自然数の二項直積]\label{k142}
$\mathbb N\times\mathbb N$ は可算無限である。
\end{thm}
\begin{proof}
$P=\mathbb N\times\mathbb N$ 上で，$(a,b)$ を $(\max\{a,b\},a,b)$ の辞書式順序によって比較する。ここで辞書式とは，最初に異なる成分の大小による比較であり，この比較は全順序となる（非反射性と全比較性は明らかであり，推移性は二つの比較で最初に異なる成分のうち先行する成分で判定すればよい）。任意の $p=(a,b)$ に対して，$p$ より小さい元の集合 $D_p$ は $S(m)\times S(m)$ に含まれる。ただし $m=\max\{a,b\}$ である。N-5 の有限集合の直積・部分集合に関する結果より $D_p$ は有限集合である。$p<q$ なら $D_p\subsetneq D_q$ であり，$|D_p|<|D_q|$ である。従って $p\mapsto |D_p|$ は $P\to\mathbb N$ の単射を定める。逆向きには $n\mapsto(n,0)$ という単射が存在するから，Cantor--Bernstein の定理により結論が従う。
\end{proof}
% genron:end K-000142
% genron:begin K-000143
% genron:use K-000001
% genron:use K-000002
% genron:use K-000003
% genron:use K-000004
% genron:use K-000005
% genron:use K-000007
% genron:use K-000017
% genron:use K-000029
% genron:use K-000033
% genron:use K-000035
% genron:use K-000039
% genron:use K-000044
% genron:use K-000045
% genron:use K-000047
% genron:use K-000048
% genron:use K-000049
% genron:use K-000100
% genron:use K-000102
% genron:use K-000107
% genron:use K-000110
% genron:use K-000118
% genron:use K-000138
% genron:use K-000141
% genron:use K-000142
\Needspace{4\baselineskip}
\begin{prop}[可算集合の有限積・有限合併]\label{k143}
有限族 $(A_i)_{i\in I}$ の各因子が高々可算ならば，その直積及び合併は高々可算である。
\end{prop}
\begin{proof}
二つの高々可算集合については，各々の $\mathbb N$ への単射を取り，直積の単射を $\mathbb N\times\mathbb N$ に入れることにより，直積が高々可算であることを得る。合併については二つの集合をそれぞれ添字 $0,1$ で区別した非交和が $\mathbb N\times\mathbb N$ に単射を持ち，合併は非交和の像であるから高々可算である。空の積は一元集合，空の合併は空集合である。有限集合 $I$ を自然数に移して帰納法を用いれば一般の場合が従う。
\end{proof}
% genron:end K-000143
% genron:begin K-000144
% genron:use K-000001
% genron:use K-000002
% genron:use K-000003
% genron:use K-000004
% genron:use K-000005
% genron:use K-000007
% genron:use K-000017
% genron:use K-000035
% genron:use K-000039
% genron:use K-000044
% genron:use K-000045
% genron:use K-000049
% genron:use K-000138
% genron:use K-000141
% genron:use K-000142
\Needspace{4\baselineskip}
\begin{prop}[指定された単射を持つ可算族の合併]\label{k144}
高々可算集合 $I$ を添字集合とする集合族 $(A_i)_{i\in I}$ を考える。各 $A_i$ から $\mathbb N$ への単射 $e_i$ が族として与えられているならば，$\bigcup_{i\in I}A_i$ は高々可算である。
\end{prop}
\begin{proof}
単射 $j\colon I\to\mathbb N$ を一つ取る。非交和 $D=\coprod_{i\in I}A_i$ に対して
\[ (i,a)\longmapsto(j(i),e_i(a)) \]
は $\mathbb N\times\mathbb N$ への単射である。従って $D$ は高々可算であり，その像である合併も高々可算である。ここでは単射の族が既に与えられていることが重要である。各 $A_i$ について単射が存在するというだけの仮定からそのような族を無断で選んではならない。
\end{proof}
% genron:end K-000144
% genron:begin K-000145
% genron:use K-000001
% genron:use K-000002
% genron:use K-000003
% genron:use K-000004
% genron:use K-000005
% genron:use K-000007
% genron:use K-000017
% genron:use K-000027
% genron:use K-000039
% genron:use K-000133
% genron:use K-000135
% genron:use K-000137
% genron:use K-000138
\Needspace{4\baselineskip}
\begin{prop}[冪集合の非可算性]\label{k145}
$\Pow(\mathbb N)$ は高々可算でない。
\end{prop}
\begin{proof}
もし $\Pow(\mathbb N)\to\mathbb N$ の単射が存在すれば，Cantor の定理における逆向きの単射と Cantor--Bernstein の定理を用いて $\mathbb N$ と $\Pow(\mathbb N)$ の全単射を得る。これは Cantor の定理に反する。
\end{proof}
% genron:end K-000145
\Needspace{4\baselineskip}
\par\medskip\noindent\textbf{参考文献}\par\smallskip
\noindent T. Jech, \textit{Set Theory}, Third Millennium Edition, Springer, 2003, Ch.~3 (Cardinal Numbers). 濃度の比較及び Cantor の定理に関する記述を参照した。本篇では一般の基数を集合として構成せず，可算族に関して単射の族の存在を暗黙に仮定しない。
\end{document}
